
\exercise
Prove that $-(-v) = v$ for every $v \in V\1$.


\exercise
Suppose $a \in \F{}\3$, $v \in V\1$, and $av = 0$. Prove that $a = 0$ or $v = 0$.


\exercise
Suppose $v, w \in V\1$. Explain why there exists a unique $x \in V$ such that $v + 3x = w$.


\exercise
The empty set is not a vector space. The empty set fails to satisfy only one of the requirements listed in the definition of a vector space (\ref{vectorspace}). Which one?


\exercise
Show that in the definition of a vector space (\ref{vectorspace}), the additive inverse condition can be replaced with the condition that
\[
0v = 0 \text{ for all } v \in V\1.
\]
Here the $0$ on the left side is the number $0$, and the $0$ on the right side is the additive identity of $V\1$.
\exercisecomment{The phrase a ``condition can be replaced'' in a definition means that the collection of objects satisfying the definition is unchanged if the original condition is replaced with the new condition.}


\exercise
Let $\infty$ and $-\infty$ denote two distinct objects, neither of which is in $\R{}$. Define an addition and scalar multiplication on
$\R{} \cup \{\infty, -\infty\}$ as you could guess from the notation. Specifically, the sum and product of two real numbers is as usual, and for
$t \in \R{}$ define
\vspace{-.08in}
\[
t \infty =
\begin{cases}
-\infty & \text{if } t < 0,\\
0 & \text{if } t = 0,\\
\infty & \text{if } t > 0,
\end{cases}
\quad\quad
t (-\infty) =
\begin{cases}
\infty & \text{if } t < 0,\\
0 & \text{if } t = 0,\\
-\infty & \text{if } t > 0,
\end{cases}
\]
and
\begin{align*}
t + \infty &= \infty + t = \infty + \infty = \infty,\\[-1bp]
\quad t + (-\infty) &= (-\infty) + t  = (-\infty) + (-\infty) = -\infty,\\[-1bp]
\infty + (-\infty) &= (-\infty) + \infty = 0.
\end{align*}
With these operations of addition and scalar multiplication, is $\R{} \cup \{\infty,-\infty\}$ a vector space over $\R{}$? Explain.


\exercise
Suppose $S$ is a nonempty set. Let $V^S$ denote the set of functions from $S$ to~$V\1$. Define a natural addition and scalar multiplication on $V^S\!$, and show that $V^S$ is a vector space with these definitions.


\exercise
Suppose $V$ is a real vector space.
\vspace{-.075in}
\begin{Itemize}
\item
The \emph{complexification} of $V\1$, denoted by $V\1_{\C{}}$, equals $V\1 \times V\1$. An element of $V\1_{\C{}}$ is an ordered pair $(u, v)$, where $u, v \in V\1$, but we write this as $u + iv$.\index{complexification!of a vector space}\sindex{VC@$V\1_{\C{}}$}
\item
Addition on $V\1_{\C{}}$ is defined by\label{1VC}
\[
(u_1 + i v_1) + (u_2 + i v_2) = (u_1 + u_2) + i (v_1 + v_2)
\]
for all $u_1, v_1, u_2, v_2 \in V\1$.
\item
Complex scalar multiplication on~$V\1_{\C{}}$ is defined by
\[
(a + bi)(u + iv) = (au - bv) + i(av + bu)
\]
for all $a, b \in \R{}$ and all $u, v \in V\1$.
\end{Itemize}
Prove that with the definitions of addition and
scalar multiplication as above, $V\1_{\C{}}$ is a complex vector space.
\exercisecomment{Think of\/ $V$ as a subset of\/ $V\1_{\C{}}$ by identifying\/ $u \in V$ with\/
$u + i0$. The construction of\/ $V\1_{\C{}}$ from\/ $V$ can then be thought of as generalizing the construction of\/ $\C{n}$ from\/ $\RR{n}$.}


